\begin{titlepage}
    \thispagestyle{empty}

    \begin{tikzpicture}[remember picture,overlay]
        \node[anchor=north west] at ([xshift=0.3cm,yshift=-0.4cm]current page.north west)
            {\includegraphics[height=2.8cm]{images/logos/ipn_logo.png}};

        \node[anchor=north east] at ([xshift=-0.3cm,yshift=-0.4cm]current page.north east)
            {\includegraphics[height=2.2cm]{images/logos/logo_escom.png}};
    \end{tikzpicture}

    \begin{center}

        {\fontsize{18}{22}\selectfont\bfseries INSTITUTO POLITÉCNICO NACIONAL\par}
        {\fontsize{16}{19}\selectfont\bfseries ESCUELA SUPERIOR DE CÓMPUTO\par}
        {\fontsize{16}{19}\selectfont\bfseries SUBDIRECCIÓN ACADÉMICA\par}

        \vspace{1.2cm}
    \end{center}

    % Fila superior
    \noindent
    \begin{minipage}[t]{0.40\textwidth}
        \raggedright
        {\fontsize{12}{14}\selectfont No de TT: 2026-B161\par}
    \end{minipage}%
    \begin{minipage}[t]{0.25\textwidth}
        \centering
        \vspace{0.55cm}
        {\fontsize{12}{14}\selectfont Documento Técnico\par}
    \end{minipage}%
    \begin{minipage}[t]{0.40\textwidth}
        \raggedleft
        {\fontsize{12}{14}\selectfont Ciudad de México, mayo de 2026\par}
    \end{minipage}

    \begin{center}
        \vspace{0.8cm}

        {\fontsize{18}{22}\selectfont\bfseries
        “Prototipo de sistema móvil de detección de insectos\\[0.15cm]
        descortezadores en árboles mediante visión por\\[0.15cm]
        computadora”\par}

        \vspace{1.3 cm}

       {\fontsize{14}{17}\selectfont\itshape Presentan\par}
       \vspace{0.2cm}
        {\fontsize{14}{17}\selectfont\bfseries Karime Calero Alvarez\textsuperscript{1}\par}
        \vspace{0.2cm}
        {\fontsize{14}{17}\selectfont\bfseries Ángel Damián Franco Rayas\textsuperscript{2}\par}
        \vspace{0.2cm}
        {\fontsize{14}{17}\selectfont\bfseries Miguel Ángel Gómez López\textsuperscript{3}\par}

        \vspace{1.3 cm}

        {\fontsize{14}{17}\selectfont\itshape Directores\par}

        \vspace{1.0cm}

        \begin{minipage}{0.42\textwidth}
            \centering
            {\fontsize{13}{16}\selectfont\bfseries\itshape Ing. Ariel López Rojas \par}
        \end{minipage}
        \hfill
        \begin{minipage}{0.42\textwidth}
            \centering
            {\fontsize{13}{16}\selectfont\bfseries\itshape Dra. Vanessa Alejandra\\Camacho Vázquez\par}
        \end{minipage}

        \vspace{1.0cm}

        {\fontsize{16}{19}\selectfont\bfseries RESUMEN\par}
    \end{center}

    \vspace{0.3cm}

    \noindent
    {\fontsize{12}{15}\selectfont
    El presente reporte describe la propuesta de un sistema orientado a la identificación de indicios de infestación por insectos descortezadores en árboles, integrado por una aplicación móvil para la captura y análisis en campo, y una plataforma web complementaria para la consulta histórica, la gestión de registros y el seguimiento de árboles analizados. A lo largo del documento se aborda el contexto del problema, su delimitación computacional, la fundamentación del enfoque de aprendizaje profundo y el análisis y diseño general de la solución propuesta.
    \par}

    \vspace{1.0cm}

    \noindent
    {\fontsize{12}{15}\selectfont
    \textbf{Palabras clave:} plagas descortezadoras de árboles, visión por computadora, áreas verdes, aprendizaje profundo, aplicación móvil.
    \par}

    \vfill

    \noindent\rule{0.75\textwidth}{0.4pt}

    {\fontsize{10}{12}\selectfont
    \textsuperscript{1} kcaleroa1900@alumno.ipn.mx \hspace{0.8cm}
    \textsuperscript{2} afrancor1900@alumno.ipn.mx \hspace{0.8cm}
    \textsuperscript{3} mgomezl1801@alumno.ipn.mx
    }
\end{titlepage}